%%%%%%%%%%%%%%%%%%%%%%%%%%%%%%%%%%%%%%%%%%%%%%%%%%%%%%%%%%%%%%%%%%%%%%%
%%%                                                                 %%%
%%%     Saint-Petersburg state University                           %%%
%%%     faculty of Applied Mathematics - Control Processes          %%%
%%%                                                                 %%%
%%%%%%%%%%%%%%%%%%%%%%%%%%%%%%%%%%%%%%%%%%%%%%%%%%%%%%%%%%%%%%%%%%%%%%%

\documentclass[a4paper]{article}
\usepackage{pmstyleeng}
\usepackage{cite}
\usepackage[T1]{fontenc}
\usepackage{lmodern}
\raggedbottom

\begin{document}
\hyphenation{SPbSU}

%%%%%%%%%%%%%%%%%%%%%%%%%%%%%%%%%%%%%%%%%%%%%%%%%%%%%%%%%%%%%%%%%%%%%%%
%%%                     UDC | AUTHOR | TITLE                        %%%
%%%%%%%%%%%%%%%%%%%%%%%%%%%%%%%%%%%%%%%%%%%%%%%%%%%%%%%%%%%%%%%%%%%%%%%
\udk{UDC 004.451.2}

%%%   1 author
%\author{Surname~N.}

%%%   2 author
\author{Shago~P.~E.}

%%% Title
\title{EEVDF Linux Process Scheduler implementation via sched\_ext}

\renewcommand{\thefootnote}{ }
{\footnotetext{{\it Shago Pavel Evgenevich} -- student, Saint-Petersburg state University;
e-mail: pavel@shago.dev, phone: +7(901)7532145}}

\maketitle

%%%%%%%%%%%%%%%%%%%%%%%%%%%%%%%%%%%%%%%%%%%%%%%%%%%%%%%%%%%%%%%%%%%%%%%
%%%%%%%                Article Recomendations                   %%%%%%%
%%%%%%%%%%%%%%%%%%%%%%%%%%%%%%%%%%%%%%%%%%%%%%%%%%%%%%%%%%%%%%%%%%%%%%%
%%%% WITHOUT recommendations
%\norec{}

%%%% Recommended by Professor ...
%\recprof{Surname~N.}

%%%% Recommended by associate Professor ...
\recdotz{Korkhov~V.~V.}

%%%% Recommended by senior lecturer ...
%\recsp{Surname~N.}


\Section{Introduction}
The research of scheduling algorithms remains a vital direction in
systems programming, since it is the scheduler that governs resource
utilization, system responsiveness, and the predictability of the operating system.
Advances in scheduler design directly improve system performance,
as evidenced by the Linux kernel's recent adoption of EEVDF over CFS,
yielding measurable latency and throughput gains across
billions of deployed systems \cite{EEVDF_intro}.

Furthermore, the emergence of sched\_ext represents a significant frontier
in Linux scheduler research. sched\_ext \cite{sched_ext_intro}
is a framework, merged into the
Linux kernel mainline as of version 6.12, that exposes a set of hooks within
the scheduling subsystem to which user-defined scheduling policies
can be attached at runtime via eBPF programs.

Rather than requiring kernel modification and recompilation,\\
sched\_ext allows custom scheduling logic to be loaded dynamically on a live system,
decoupling policy from the kernel's core mechanism.
This enables rapid prototyping of novel scheduling strategies.
A notable predecessor to this approach is ghOSt \cite{ghOSt}, a framework
developed at Google that delegates scheduling decisions to userspace agents.

Conducted research combines reimplementation of the
Earliest Eligible Virtual Deadline First (EEVDF) algorithm, the Linux fair-class scheduler
since version 6.6 \cite{EEVDF_intro}, through the sched\_ext framework.


\Section{Model}
EEVDF, formalized by Stoica and Abdel-Wahab in~\cite{stoica, stoica2}, is a proportional-share
discipline for a time-shared resource dispatched in quanta of size at most $ q $.
Each active client $ i $ carries a positive weight $ w_i $, and the aggregate
active weight at time~$ t $ is
\begin{equation}
    W(t)=\sum_{j\in A(t)} w_j ,
\end{equation}
where $ A(t) $ denotes the set of clients competing for service.
The ideal reference is a continuous-flow system in which client
$ i $ continuously receives instantaneous share $ w_i/W(t) $.
EEVDF captures this entitlement through system virtual time
\begin{equation}
    V(t)=\int_{t_0}^{t}\frac{1}{W(\tau)}\,d\tau ,
\end{equation}
so that one unit of $ V $ corresponds to one unit of service per unit weight.

Let $ t_0^i $ denote the instant at which client $ i $ most recently became
active and $ s_i(t_0^i,t) $ the actual service it has accumulated since that
instant. When client $ i $ submits a request of size $ r $ at time $ t $,
the algorithm assigns a virtual eligible time
\begin{equation}
    v_e = V(t_0^i)+\frac{s_i(t_0^i,t)}{w_i}
\end{equation}
and a virtual deadline
\begin{equation}
    v_d = v_e + \frac{r}{w_i} .
\end{equation}
The request is considered runnable only once $ v_e \le V(t) $.
Among all eligible requests, EEVDF dispatches the one with the
smallest~$ v_d $. This rule couples fairness with latency.
Clients that have received less than their continuous model share become
eligible sooner, while larger weights compress the term $ r/w_i $,
producing earlier deadlines for the same request size.
The lag of client~$ i $, defined as the signed deviation between ideal
and actual service, satisfies
\begin{equation}
    |\mathrm{lag}_i(t)|<q
\end{equation}
in steady state, an optimal bound for any discrete proportional-share
scheduler operating with quantum~$ q $~\cite{stoica, stoica2}.

\enlargethispage{-1\baselineskip}
\Section{Mapping}
The proposed sched\_ext variant implements EEVDF as a discrete, event-driven scheduler
whose state consists of a global virtual-time reference and per-task virtual
service timestamps.

In the implementation, the client is a runnable task,
$ w_i $ is the task's scheduler weight (kernel nice value), and its virtual
position is stored in thread-local \verb|scx.dsq_vtime|.
Global state is minimal and consists only of \verb|vtime_now|, which plays
the role of the system virtual-time frontier, and \verb|total_weight|, which 
tracks the total runnable weight.

The theoretical request size $ r $ is represented by the execution slice
$ \Delta $ assigned to a task by sched\_ext subsystem.
On enqueue, the scheduler computes a discrete analogue of the EEVDF eligibility
condition,
\begin{equation}
    VE_i = \max\bigl(v_i,\; V_{\mathrm{now}} - \Delta\bigr),
\end{equation}
where $ v_i $ is the task's stored virtual time and
$ V_{\mathrm{now}} \equiv \verb|vtime_now| $.
This prevents a task from re-entering the system with an excessively old
virtual timestamp while preserving the effect of lag.
The corresponding virtual deadline is then
\begin{equation}
    VD_i = VE_i + \frac{\Delta \cdot S}{w_i},
\end{equation}
where $ S $ is an integer scaling constant used in place of fractional
arithmetic.
As in the analytical EEVDF rule, increasing $ w_i $ reduces the deadline
increment, so higher-weight tasks return earlier in virtual order and obtain a
proportionally larger share of processor time.

Task selection is realized through sched\_ext dispatch queues.
The scheduler maintains a single shared DSQ ordered by $ VD_i $ via\\
\verb|scx_bpf_dsq_insert_vtime()|, so dispatch from that queue implements
earliest virtual deadline first.
If \verb|select_cpu()| finds an idle CPU, the task may instead be inserted
directly into the local DSQ, reducing wakeup latency without changing the
underlying policy.

Consumed service is charged in \verb|stopping()|:
\begin{equation}
    v_i \leftarrow v_i + \frac{\mathrm{consumed}\cdot S}{w_i}.
\end{equation}
This is the core fairness mechanism, so heavier tasks accumulate virtual time more
slowly for the same physical runtime.
Unlike the continuous model, \verb|vtime_now| is not maintained by explicit
integration of $ 1/W(t) $. It is advanced lazily in \verb|running()| to at least
the selected task's virtual time.

\enlargethispage{-1\baselineskip}
\Section{Results}
In this section, three benchmarks are selected for analysis:
hackbench (1000 tasks), sysbench (default parameters) and sched-latency.

Hackbench is a well-known Linux scheduler benchmark designed to
measure scheduler and inter-process communication (IPC) performance.
It creates groups of processes that exchange messages intensively, simulating high system
load. The total execution time reflects how efficiently the scheduler handles
context switching and task coordination.

Sysbench is a versatile benchmarking tool used to evaluate system performance
across CPU, memory, and disk I/O. It generates computational workloads
that simulate real-world usage scenarios, such as
online transaction processing (OLTP) database operations, making it well
suited for assessing overall system throughput under realistic conditions.

sched-latency is a self-written benchmark that uses eBPF tracing capabilities
to measure scheduling latency. It attaches to kernel tracepoints (\verb|sched_wakeup|, \verb|sched_switch|) and enqueue hooks to record the schedule delay for both
the implemented sched\_ext EEVDF scheduler and the baseline kernel EEVDF.

\Figure{0.9\textwidth}{images/Fig1.eps}
{Test results for builtin Linux 7.0-rc4 EEVDF scheduler and
implemented sched\_ext EEVDF}


\Section{Conclusion}
The proposed sched\_ext-based EEVDF implementation shows no measurable
degradation in hackbench (0.20\,s for both schedulers, +0.0\%), demonstrating
that the eBPF dispatch overhead is negligible for intensive, scheduler demanding workloads.

Sysbench throughput increased from 4157 events/s (default EEVDF) to
7357 events/s (sched\_ext EEVDF), a +77\% improvement, suggesting that
the sched\_ext variant may offer advantages for sustained computational workloads.

The 99th percentile scheduling latency of the sched\_ext variant is
approximately 250\,$ \mu $s compared to 25\,$ \mu $s for the in-kernel
baseline, a 10$ \times $ increase attributable to the overhead
of eBPF dispatch and contention on the single shared dispatch queue (DSQ).

However, the main achievement of this research and proposed sched\_ext EEVDF
is the flexibility of the resulting solution. Any constant (e.g.\ scaling factor $ S $)
or formula (e.g.\ deadline rule $ VD_i = VE_i + \Delta S / w_i $) can be modified and tested
without requiring a reboot, thanks to the sched\_ext framework.


%%%% References
\begin{thebibliography}{7}

\bibitem{EEVDF_intro}
An EEVDF CPU scheduler for Linux [Internet resource]:\\
\url{https://lwn.net/Articles/925371} (date 10.11.2025).

\bibitem{sched_ext_intro}
The extensible scheduler class
[Internet resource]:\\
\url{https://lwn.net/Articles/922405} (date 03.09.2025).

\bibitem{ghOSt}
Humphries~J.~T., Natu~N., Chaugule~A., Weisse~O., Rhoden~B., Don~J., Rizzo~L., Rombakh~O., Turner~P., Kozyrakis~C. ghOSt: Fast \& Flexible User-Space Delegation of Linux Scheduling // Proceedings of the ACM SIGOPS 28th Symposium on Operating Systems Principles (SOSP '21). 2021. P. 588-604.

\bibitem{stoica} Stoica~I., Abdel-Wahab~H.
Earliest eligible virtual deadline first: A flexible and accurate mechanism
for proportional share resource allocation. PhD thesis. USA: Old Dominion University, 1995. P. 1-37.

\bibitem{stoica2}
Stoica~I., Abdel-Wahab~H., Jeffay~K., Baruah~S., Gehrke~J., Plaxton~C. G. A Proportional Share Resource Allocation Algorithm for Real-Time, Time-Shared Systems // Proceedings of the 17th IEEE Real-Time Systems Symposium. 1996. P. 1-12.

%\bibitem{hackbench}
%Hackbench: A new multiqueue scheduler benchmark
%[Internet resource]:\\
%\url{http://www.lkml.org/archive/2001/12/11/19/index.html} (date 17.11.2025).


\end{thebibliography}


\end{document}
